\chapter{The corpus}

This chapter states what each document in the corpus is, as established by opening it, and separates
that from what its filename says. The corpus is seven files, 7{,}852{,}267 bytes.

\section{What is in it}

Each document is identified from its own first page and embedded metadata instead of its filename.
Five are the official Clay Mathematics Institute problem
descriptions~\cite{src:Cook,src:Deligne,src:Bombieri,src:Jaffe-and-Witten,src:Wiles}. Two are
Riemann~\cite{src:Wilkins-1998,src:Riemann-1859}.

\begin{center}
\begin{tabular}{@{}llrr@{}}
\toprule
File & Identified as & Pages & Bytes \\
\midrule
\texttt{pvsnp.pdf}       & Stephen Cook, \emph{The P versus NP Problem}      & 12 & 176{,}204 \\
\texttt{hodge.pdf}       & Pierre Deligne, \emph{The Hodge Conjecture}       & 5  & 146{,}894 \\
\texttt{riemann.pdf}     & E. Bombieri, \emph{Problems of the Millennium:}   & 11 & 159{,}267 \\
                         & \emph{the Riemann Hypothesis}                     &    &           \\
\texttt{yangmills.pdf}   & Jaffe and Witten, \emph{Quantum Yang--Mills}      & 14 & 195{,}976 \\
                         & \emph{Theory}                                     &    &           \\
\texttt{birchswin.pdf}   & Andrew Wiles, \emph{The Birch and}                & 5  & 142{,}908 \\
                         & \emph{Swinnerton-Dyer Conjecture}                 &    &           \\
\texttt{Wilkins-}        & Riemann 1859, \emph{On the Number of Prime}       & 10 & 119{,}504 \\
\texttt{translation.pdf} & \emph{Numbers less than a Given Quantity},         &    &           \\
                         & translated by David R. Wilkins                    &    &           \\
\texttt{riemann1859.pdf} & Not the 1859 paper. See below.                    & 4  & 6{,}911{,}514 \\
\bottomrule
\end{tabular}
\end{center}

\section{The corpus is complete}

Six Millennium problems are commonly counted as open, and the corpus holds five official statements.
The natural first reading is that Navier--Stokes is missing. That reading is wrong.

The Clay Mathematics Institute index of Millennium Prize problems~\cite{src:Clay-Mathematics-Institute},
accessed 11 September 2026, lists five problems as unsolved, Birch and Swinnerton-Dyer, Hodge, P
versus NP, Riemann, and Yang--Mills and the mass gap, and one as solved, Poincar\'e. Navier--Stokes
appears in neither list. The corpus holds exactly the five the index lists as unsolved.

That is a fact about how a website categorizes its own contents on one date at one address, and it is
recorded as that and not as a fact about the status of a problem. A site can be mid-edit,
deliberately cautious, or behind. The Navier--Stokes problem's own page, accessed the same day,
carries the status label \emph{Active}. The two readings are not consistent with each other, both
are recorded, and neither is treated here as evidence about the mathematics.

\section{The document that is not what it is named}

\texttt{riemann1859.pdf} is not Riemann's 1859 memoir. It is a photographic scan of Riemann's
handwritten working manuscript~\cite{src:Riemann-1859}.

The evidence is in the file. Its four pages carry zero extractable characters between them, measured
page by page. The embedded document title is a scanner path,
\texttt{Cod\_\allowbreak{}Ms\_\allowbreak{}Riemann\_\allowbreak{}3\_\allowbreak{}jpg}, under a
directory named for the Riemann \emph{Nachlass} and a subdirectory named \emph{Rechenseiten},
calculation pages. The producer chain is a JPEG passed through PDF24 Creator and a Quartz PDF context
in 2015, which is a scan-to-PDF path and not a typesetting one. Rendered and read as images, all four
pages are German cursive in ink, heavily struck through with interlinear corrections, foxed and
water-stained, with \emph{Nachlass Riemann} and the year 1859 legible in the header of the first
page. The zeta function, the factor $\Pi(s-1)$, the logarithmic expansion of the Euler product, the
function $\xi(t)$ and the logarithmic integral $\mathrm{Li}(x)$ are identifiable in the hand.

The file is a primary source in the most literal available sense, and it is unusable for text search
or for quotation without a reader of nineteenth-century German \emph{Kurrentschrift}. The content of
the published memoir is in the Wilkins translation~\cite{src:Wilkins-1998} instead.

\textbf{A citation built from this filename is wrong.} It would name Riemann's published memoir and
point at a manuscript draft. A bibliographic entry left with empty fields until somebody looks the
source up does not catch the error, because the error is in the identification and not in the
completeness. The distinction has to be written into the entry, and it is in this paper's.

\section{An observation about production}

Five of the six official statements, counting the Navier--Stokes statement~\cite{src:Fefferman}
fetched separately for the next chapter, share a producer string of MiKTeX pdfTeX-1.21a and creation
timestamps on 2026-08-04 within three minutes of one another; the Navier--Stokes file is stamped
16:46:57, nineteen seconds after the Birch and Swinnerton-Dyer file. Bombieri's Riemann statement
comes from older tooling, a dvips and Acrobat Distiller chain dated 2000-08-30, and carries the
embedded title \texttt{Riemann-new.dvi}.

The observation is about how the statements are produced and says nothing about their contents. It
establishes that the corpus and the separately fetched statement are the same edition of the same
set, which matters for a bibliographic entry and for no argument about mathematics.
